% Decisiones de diseño e integración (equipo).
\section{Decisiones de diseño}

Este capítulo reúne las decisiones que el equipo tomó al traducir los modelos y
al evaluar las restricciones de integridad. Ninguna de ellas se sigue de manera
automática del enunciado: todas proceden de una lectura deliberada de los datos
disponibles y deben declararse para que el reporte sea reproducible. Cada
módulo documenta aquí sus propias decisiones; las convenciones generales y el
estado de avance se registran al inicio.

\subsection{Convenciones generales del reporte}

\begin{itemize}
  \item \textbf{Entidades y relaciones.} Los nombres de entidad se escriben en
        singular y sin artículo, de modo que la entidad \emph{Departamento} no
        se confunda con una instancia. Las relaciones que expresan una acción
        se nombran con el verbo en infinitivo, como en \emph{Trabajar},
        \emph{Administrar} y \emph{Estudiar}.
  \item \textbf{Atributos.} Se escriben en una sola palabra, con inicial en
        mayúscula y sin espacios ni guiones, del tipo \emph{NombreDepto} o
        \emph{FechaContratacion}. Las abreviaturas se conservan sin punto.
  \item \textbf{Llaves.} La llave primaria se marca de forma explícita en el
        diagrama y se declara en el esquema; los atributos que la componen
        conservan el nombre del diagrama de origen. Las llaves foráneas
        coinciden con la llave primaria que referencian y se identifican en el
        diagrama con la notación de llave foránea.
  \item \textbf{Atributos de tiempo.} Los periodos se representan con el par
        \emph{Desde} y \emph{Hasta}, con \emph{Desde} inclusivo y \emph{Hasta}
        exclusivo, según la convención $[desde, hasta)$ fijada en el enunciado.
\end{itemize}

\begin{conceptobox}[Divergencia de nomenclatura entre entregas]
Los diagramas de la Tarea~02 escriben sus atributos sin acentos, con
\emph{NumProf} y \emph{Titulo} como ejemplos, mientras que la figura del
Ejercicio~4 de esta tarea los emplea con acento, con \emph{NúmEmpleado} y
\emph{Género}. Ambas formas aparecen en el material del equipo y no es posible
unificarlas sin modificar una entrega anterior ya evaluada. Se adopta el
criterio de conservar la grafía de la fuente: cada módulo emplea literalmente
los nombres del enunciado o del diagrama de Draw.io que le corresponde.
Resolver esta divergencia exige una decisión colectiva en la próxima entrega,
que debería fijar una sola política de acentos en los nombres de atributo.
\end{conceptobox}

\subsection{Estado de los módulos}

El reporte integra las respuestas de los cinco integrantes y los diagramas
disponibles para los ejercicios de traducción.

\subsection{Módulo 1. Ejercicio~1 y Ejercicio~4, incisos a a e}

\begin{itemize}
  \item \textbf{Lectura de la notación de cardinalidad.} La figura del
        Ejercicio~4 emplea dos marcadores: el par $\circ|$ junto a la entidad
        del lado uno, que indica cardinalidad máxima uno, y la pata de cuervo
        junto a la entidad del lado muchos. El círculo exterior se interpretó
        como participación parcial. Dado que \emph{NúmEmpleado} y
        \emph{NúmDepto} forman parte de las llaves primarias, las llaves
        foráneas se declararon obligatorias. Se dejó constancia de que los
        veredictos de los incisos~b y~d no dependen de esta lectura.
  \item \textbf{El atributo \emph{Hasta} fuera de la llave primaria.} En
        \emph{Salario}, \emph{Trabajar} y \emph{Administrar}, el atributo
        \emph{Hasta} no forma parte de la llave primaria, tal como aparece en
        la figura del enunciado. La decisión tiene una consecuencia que el
        ejercicio hace visible: el esquema admite periodos solapados siempre
        que las fechas de inicio sean distintas.
  \item \textbf{Atributo no clave.} En \emph{Departamento}, el atributo
        \emph{NombreDepto} no se declaró único. Es la decisión que permite el
        inciso~a y se documentó como tal, junto con la restricción adicional
        que la haría imposible.
  \item \textbf{Llaves candidatas no declaradas.} En \emph{Empleado} se
        documentó que la combinación de nombres y fecha de nacimiento es una
        llave candidata que el esquema no elige como primaria, en lugar de
        omitir la observación.
  \item \textbf{Restricción no modelada en el inciso~d.} Se documentó que la
        exigencia de que quien administra un departamento trabaje en él no
        puede expresarse como llave foránea, porque el par
        (\emph{NúmEmpleado}, \emph{NúmDepto}) es superllave de \emph{Trabajar}
        pero no llave candidata al incluir \emph{Desde} en su llave primaria.
        Se propuso como alternativa una aserción.
\end{itemize}

\subsection{Módulo 2. Ejercicio~2.a}

Las reglas aplicadas y las decisiones de traducción se presentan en el
Ejercicio~2.a, junto con el diagrama relacional.

\subsection{Módulo 3. Ejercicio~2.b, Sistema de Información Geográfica}

Las reglas aplicadas y el tratamiento de la relación recursiva Alimentar
se presentan en el Ejercicio~2.b, junto con los diagramas del SIG.

\subsection{Módulo 4. Ejercicio~3, incisos a a c}

\begin{itemize}
  \item \textbf{Llave compuesta de \emph{A}.} El atributo subrayado $a$ del
        modelo E/R se descompuso en $(a_1,a_2)$; la pareja completa identifica
        a \emph{A} y se conserva como llave foránea compuesta donde procede.
  \item \textbf{Orden de atributos.} Se siguió el formato del enunciado:
        llaves primarias, atributos propios y llaves foráneas. Así,
        $\mathrm{AB}(a_1,a_2,b,ab_1)$ en $M:N$ y
        $\mathrm{B}(b,b_1,ab_1,a_1,a_2)$ en $1:N$ fijan la posición de cada
        valor en los incisos~b y~c.
  \item \textbf{Participación parcial.} Para $1:N$ se permite que la pareja
        foránea $(a_1,a_2)$ esté ausente de \emph{B}, pero no que se almacene
        solo uno de sus componentes. En $1:1$ se eligió una relación
        \emph{AB} independiente, con llave primaria $(a_1,a_2)$ y unicidad
        sobre $b$; la ausencia de fila representa la no participación.
  \item \textbf{Estado inicial del inciso~c.} Se conservó \emph{A} del
        inciso~b para comprobar las llaves foráneas, pero se evaluó cada
        conjunto sobre una \emph{B} vacía del nuevo modelo $1:N$. Las cinco
        tuplas de \emph{B} citadas en el inciso~b corresponden al escenario
        $M:N$ y no se arrastraron como inserciones previas.
\end{itemize}

\subsection*{Módulo 5. Ejercicio 3, inciso d}

Para las 5 tuplas válidas e inválidas, se tomó en cuenta el escenario
dado por el inciso b), donde las llaves de las tuplas de
\textit{A} son $(2,\text{'ww'})$, $(4,\text{'xx'})$,
$(6,\text{'yy'})$ y $(8,\text{'zz'})$, mientras que las llaves de
las tuplas de \textit{B} son $17$, $27$, $37$, $47$ y $57$;
además se tomó como base el modelo relacional $1:1$ y parcial en
ambos sentidos, por lo cual los elementos de \textit{A} y
\textit{B} pueden o no participar en la relación \textit{ab} y si
participan una tupla de \textit{A} no puede estar relacionada con
más de una tupla de \textit{B} y viceversa una tupla de \textit{B}
no puede estar relacionada con más de una tupla de \textit{A}.

\subsection*{Módulo 7. Ejercicio\textasciitilde 4, incisos f a j}

Se tomó al atributo \textit{Hasta} de las tablas \textit{Salario},
\textit{Trabajar} y \textit{Administrar}, como una fecha sin
restricciones visibles en el modelo relacional presentado.
